\section{The same common leading functions give the whole weighted slope inequality}

Let $C$ be an original integral normal projective curve over an arbitrary
field $k$, with its actual structure map $p:C\to\operatorname{Spec}k$.
Normality means that every original stalk is integrally closed. The
projective factorization of this SAME $p$ derives properness, finite type,
local Noetherianity and compactness. In particular the SAME original $C$
is Noetherian. The dimension bound is $\dim C\le1$; the dimension-one
case is the intended curve case.

\begin{definition}[The whole original residue-degree weight]
\label{def:gauss-common-leading-original-residue-weight}
\uses{def:fc-affine-scheme-local-coefficients,def:fc-original-weighted-rational-order}
\lean{EnvelopeProperty.GaussCommonLeadingDegree.residueWeight}
\leanok
For every original point $x\in C$ put
\[
 w_p(x)=\dim_k\kappa_C(x),
\]
using the SAME original coefficient action induced by the actual $p$.
In real-valued sums this is the real cast of the literal native finite
rank. It is nonnegative and is not replaced by $1$. At closed points
these are the actual residue degrees over $k$; the generic-point term
below vanishes because its native order is zero. No residue-field
isomorphism, rational-point or unramified-branch hypothesis is supplied.
For every original rational function $g$,
\[
 w_p(x)\operatorname{ord}_x(g)
   =\bigl(\operatorname{weightedOrder}_{p,g}(x):\mathbf R\bigr).
\]
\end{definition}
\begin{proof}
Use the existing actual residue algebra of $p$ at the SAME original
point. Its coefficient is the original evaluation of the original base
scalar. Take the native rank over this exact action and cast it to the
real numbers. Nonnegativity is the nonnegativity of a natural-number
cast. The displayed identity follows by casting the definition of the
old original integer weighted order and commuting the two real factors.
Every original point and its original field are retained.
\end{proof}

\begin{theorem}[True finite support and real principal balance on the same curve]
\label{thm:gauss-common-leading-real-principal-balance}
\uses{thm:fc-original-curve-principal-weighted-balance,def:fc-original-weighted-rational-order,def:gauss-common-leading-original-residue-weight}
\lean{EnvelopeProperty.GaussCommonLeadingDegree.principal_balance_real}
\leanok
For EVERY nonzero original $g\in C.\mathrm{functionField}$, the actual
function $x\mapsto\operatorname{ord}_x(g)$ has finite support and
\[
 \sum_{x\in C}w_p(x)\operatorname{ord}_x(g)=0
\]
as a real finite-support sum. Neither support finiteness, a Cartier
denominator nor a principal-divisor degree balance is an input.
\end{theorem}
\begin{proof}
The actual projective factorization of $p$ derives the Noetherian
hypotheses on $C$. Apply the already proved original rational-order
finite-support theorem, which internally produces the two Cartier
denominators from the actual normal-curve geometry.
Apply the existing full original weighted principal-divisor theorem to
this SAME nonzero $g$. That theorem derives the two genuine Cartier
degree sums from the isomorphism of the original reciprocal denominator
modules and their FULL original Euler characteristics. It already gives
zero for the native integer weighted sum; no such degree argument is
reproved here.

Take the finite set given by the actual unweighted order support. Outside
it both the native integer weighted function and the real weighted
function vanish. Replace both whole original sums by their sums over
that SAME finite set. The preceding original coefficient identity and
finite-sum cast give the real cast of the old integer sum, hence zero.
This retains every actual original residue-field degree.
\end{proof}

\begin{theorem}[Two true principal balances control the entire active slope]
\label{thm:gauss-common-leading-global-weighted-max}
\uses{thm:gauss-common-leading-real-principal-balance,def:gauss-common-leading-original-residue-weight,analysis:actual-finite-affine-active-max,analysis:actual-weighted-active-max-slope-sum,thm:fc-curve-proper-support-points}
\lean{EnvelopeProperty.GaussCommonLeadingDegree.weighted_activeSlope_nonneg}
\leanok
Let $B_1,\ldots,B_n,B_t\ne0$ be actual original functions in
$C.\mathrm{functionField}$ and let $r_1,\ldots,r_n\in\mathbf R$ be
fixed BEFORE $x$. Let $I\subseteq\mathbf Q^n\times\mathbf Q$ be finite
and nonempty and let $\beta:(\mathbf Q^n\times\mathbf Q)\to\mathbf R$
be one fixed original intercept family. Put
\[
 \begin{split}
  \gamma_{\ell,x}&=\operatorname{ord}_x(B_\ell)
                   -r_\ell\operatorname{ord}_x(B_t),\\
  s_j(x)&=-\sum_\ell j_\ell\gamma_{\ell,x},\\
  I_0&=\{j\in I:\beta_j=\max_{i\in I}\beta_i\},\qquad
  D(x)=\max_{j\in I_0}s_j(x).
 \end{split}
\]
These are the actual old active set and active maximum. Then the actual
$D:C\to\mathbf R$ has finite support, $D(\eta_C)=0$, and
\[
 \sum_{x\in C}w_p(x)D(x)\ge0.
\]
The sum is over ALL original points, equivalently all nongeneric closed
points because the generic term is zero. Negative original orders and
arbitrary fixed $r_\ell$ are included. The case $n=0$ is included.
No individual zero balance or finite-support hypothesis is assumed.
\end{theorem}
\begin{proof}
Apply the preceding true principal theorem to EACH $B_\ell$ and to the
SAME common $B_t$. Let $S$ be the union of their actual finite native-order
supports. It is finite even when $n=0$. Outside $S$ every original order
in the displayed formulas is zero, hence every $\gamma_{\ell,x}$ and
every $s_j(x)$ is zero. The actual active set is nonempty. Its maximum
of the constant zero family is zero, so $D(x)=0$ outside $S$.
Thus the actual support of $D$ is a subset of this true finite set.
The generic point is the top point in the original specialization order
and has coheight zero. The native integer order is zero off coheight
one. Therefore all the same orders, slopes and their active maximum are
zero at the actual generic point.

All actual weighted functions involved vanish outside $S$. The TWO
true principal balances for $B_\ell$ and $B_t$, on this SAME finite set,
give for every $\ell$
\[
 \sum_{x\in S}w_p(x)\gamma_{\ell,x}
 =\sum_{x\in S}w_p(x)\operatorname{ord}_x(B_\ell)
    -r_\ell\sum_{x\in S}w_p(x)\operatorname{ord}_x(B_t)=0.
\]
Every $r_\ell$ is constant in $x$. Distribute the finite rational linear
combination and interchange the two finite sums. For EVERY original $j$
this derives, rather than assumes,
\[
 \sum_{x\in S}w_p(x)s_j(x)
   =-\sum_\ell j_\ell\sum_{x\in S}w_p(x)\gamma_{\ell,x}=0.
\]
Apply the EXISTING weighted active-maximum theorem to the SAME actual
$I,\beta,S,s,w_p$. Its weight nonnegativity is the original rank
nonnegativity and its individual balances are the just-derived equalities.
It gives $\sum_{x\in S}w_p(x)D(x)\ge0$. Outside $S$ the latter original
weighted function vanishes, so its whole finite-support sum is this
SAME finite sum. This proves the claim, without reproving the finite
maximum inequality. For $n=0$ all coordinate sums are empty and the same
argument gives zero slopes everywhere.
\end{proof}

\begin{theorem}[The original common Gauss profiles have nonnegative whole residue-weighted degree]
\label{thm:gauss-common-leading-original-degree-nonnegative}
\uses{thm:gauss-common-global-leading-convex-profiles,thm:gauss-common-leading-global-weighted-max,def:gauss-common-leading-original-residue-weight,scheme:native-regular-unit-local-psh,analysis:actual-finite-affine-active-max}
\lean{EnvelopeProperty.GaussCommonLeadingDegree.exists_common_leading_degree_nonneg}
\leanok
Keep the actual original integral proper dominant
$f:Y\to\operatorname{Spec}k[t]$, actual integral closed immersion
$i:C\hookrightarrow Y$, true vertical $i(\eta_C)$, and genuine original
local component geometry used by the frozen common-leading profile
producer. Assume the SAME original $C$ is normal of dimension at most
one and is genuinely projective over $k$ for its LITERAL structure map
\[
 p_C=C\xrightarrow{i}Y\xrightarrow{f}\operatorname{Spec}k[t]
                \longrightarrow\operatorname{Spec}k.
\]
Let $u$ have the original native convex-PL property on $U$ at the prescribed
paper Gauss radius. The old actual common $q_C$ is retained. Under its
honest original applicability condition $q_C\in U$, retain the SAME
$W,V$, original unit family, rational $I$, intercepts, active set,
global $\Pi,e>0,U_t,A_\ell,U_\ell$, and original native profiles.

The actual leading functions are the ORIGINAL generic stalk images
$B_\ell=i^\#_{\eta_C}(U_\ell)$ and $B_t=i^\#_{\eta_C}(U_t)$.
They are nonzero actual elements of the SAME $C.\mathrm{functionField}$.
For the SAME actual local slopes
\[
 D_x=\max_{j\in I_0}
       \left(-\sum_\ell j_\ell\left(
        \operatorname{ord}_x(B_\ell)
          -\frac{A_\ell}{e}\operatorname{ord}_x(B_t)\right)\right),
\]
internally derive finite support on the WHOLE original curve and
\[
 D_{\eta_C}=0,\qquad
 \sum_{x\in C}[\kappa_C(x):k]D_x\ge0.
\]
The actual original weighted points, scales, values of $u$, exact
quotients and their limits remain those of the already proved common
profile theorem, with slope EXACTLY this $D_x$. There is no added
support, balancing, global-leading matching, degree or norm certificate.
All original closed-point residue degrees are retained.
\end{theorem}
\begin{proof}
The genuine projective factorization of the literal $p_C$ supplies the
same $C$ with its required local Noetherianity and compactness.
Call the frozen common-leading profile producer once with the actual
original geometry and original convex-PL hypothesis. It supplies the
actual $q_C$. Under its original applicability condition it supplies
all the SAME global and local data in the statement. In particular its
true generic-stalk unit factors, global exponents and positive $e$ are
chosen before every original local point.

The actual generic closed-immersion stalk map takes each original unit
to a unit of the actual original generic stalk of integral $C$. Thus
every $B_\ell$ and the SAME $B_t$ is nonzero. Apply the preceding actual
weighted-maximum theorem to THESE functions, the fixed real coefficients
$r_\ell=A_\ell/e$, THESE actual original intercepts and active indices,
and the literal original $p_C$. It derives the claimed whole finite
support, generic zero and residue-weighted sum inequality internally.
No vertical exponent is set to zero, and no original residue weight is
replaced by a rational-point or degree-one weight.

The active slope there has direction set the whole original $C$.
At a fixed original point, evaluation of this same finite maximum is
literally the active slope with singleton direction parameter used by
the frozen profile theorem. Substitute this definitional equality in
the old exact original function profiles, actual quotients and limits.
All points, functions, scales, indices and original whole-field
factorizations remain the same. This supplies precisely the claimed
global weighted slope consumer.
\end{proof}

This chapter does not identify a generic section with a model-line
frame, compare its orders with a line, or compute the degree of an
original line bundle. Those separate dictionary and degree bridges
are needed before a residual line-degree or LocalPositivity conclusion.
